% 本文件由 LaTeX 排版 Agent 根据有效 translation.json 自动生成。
% author=Lactantius; work_id=0703; title=《论天主的忿怒》

\chapter{《论天主的忿怒》}

% structure: p0001 source=h1 target=omitted reason=page-title-already-rendered-by-parent

致多纳图斯\par

% structure: p0003 source=h2 target=section reason=relative-heading-rank; base=h2

\section{第一章——论神圣智慧与人的智慧。}

多纳图斯，我常观察到，许多人持有这一见解，某些哲学家也曾主张，即天主不会发怒；因为天主性是全然慈善的，并且与祂那卓越超绝的权能相悖的是，伤害任何人；或者至少，祂完全不留意我们，以至于我们从祂的良善中得不到任何益处，也从祂的恶意中得不到任何灾祸。然而，这些人的错误极大，且倾向于颠覆人类生活的状态，我们必须予以驳斥，免得你也被那些自认为智慧的人的权威所煽动而受骗。不过，我们并非如此傲慢，竟夸耀靠自己的理智就领悟了真理；我们乃是追随天主的教导，唯有祂能够知晓并启示隐秘之事。但哲学家们由于缺乏这一教导，便臆测事物的本性可由推测来确定。但这是不可能的；因为人的心智，被囚禁在身体的黑暗寓所中，远远不能感知真理；在此，天主性与人性有别：无知属于人性，知识属于天主性。\par

因此，我们需要某种光明，以驱散笼罩着人思虑的黑暗，因为当我们生活在有死的肉身中时，无法凭感官洞察。然而，人灵的光明就是天主；谁认识祂并将祂迎入心中，谁就能以明亮的心承认真理的奥秘；但当天主和神圣的训导被移除后，一切便充满了错误。\Person{苏格拉底}虽然是所有哲学家中最博学的，但为了证明其他那些自认为有所知的人的愚昧，他恰当地说他什么都不知道，只知道一件事——即他什么都不知道。因为他明白，那学问本身没有确定、没有真实的东西；他并非像某些人想象的假装学问以驳斥他人，而是他某种程度上看到了真理。他甚至在其审判中（如\Person{柏拉图}所述）作证说，没有人的智慧。他如此鄙视、嘲笑并抛弃了当时哲学家们所夸耀的学问，以至于将此事本身称为最大的学问，即他学会了知道自己一无所知。因此，如果像\Person{苏格拉底}所教导、\Person{柏拉图}所传授的那样，没有人的智慧，那么显然，真理的知识是神圣的，只属于天主，不属于任何其他。因此，必须认识天主，唯有祂内才有真理。祂是世界的父母，万物的创造者；祂不能被眼睛看见，也难以被心智分辨；那些既未曾获得真智慧，又未能理解伟大而神圣奥秘体系的人，惯常以多种方式攻击祂的宗教。\par

% structure: p0006 source=h2 target=section reason=relative-heading-rank; base=h2

\section{第二章——论真理及其阶梯，并论天主。}

因为既然有许多阶梯，藉以攀登至真理的住所，任何人都不易达到顶峰。当眼睛被真理的光辉所眩惑时，那些无法保持稳定脚步的人便跌回到平地上。现在，第一步是理解虚假的宗教，并抛弃偶像崇拜，即人手所造的神明的邪神崇拜。第二步是以心智洞察只有一位至高的天主，祂的权能和眷顾从起初创造了世界，并随后继续治理它。第三步是认识祂的仆人和使者，即祂派遣到大地作为祂的使节的那一位，藉着祂的教导，我们得以摆脱曾纠缠我们的错误，并受教于真天主的崇拜，从而学会正义。正如我所说，从所有这些阶梯，若不将双脚以坚定不移的稳固牢牢扎稳，便极易滑倒堕落。\par

我们看到那些从第一步上跌落的人，他们虽然理解事物是虚假的，却并未发现真实的事物；他们虽鄙视尘世和易朽的偶像，却未投身于他们所不认识的天主的崇拜。反而惊叹于宇宙的元素，崇拜天、地、海、太阳、月亮及其他天体。\par

然而，我们已在\WorkTitle{神圣训导}第二卷中谴责了他们的无知。但我们说，那些从第二步上跌落的人，虽明白只有一位至高的天主，却被哲学家所诱骗，被虚假的论据所迷惑，对那崇高威严持有远非真理的看法；他们要么否认天主有任何形状，要么认为祂不受任何情感影响，因为每种情感都是软弱的标志，而在天主内并无软弱。但从第三步上跌落的人，则是那些认识天主的使者，即那神圣不朽殿堂的建造者，却或不接受祂，或以与信德要求不同的方式接受祂；我们在上述著作的第四卷中已部分地驳斥了他们。将来，当我们开始答复所有在争论中毁灭了真理的派别时，我们将更加仔细地予以驳斥。\par

但如今我们要反驳那些从第二步跌倒，对至高天主怀有错误看法的人。因为有些人说祂既不向任何人施恩，也不发怒，而是在安稳与宁静中享受自己不朽的益处。另一些人虽除去义怒，却将慈爱归于天主；他们认为，一个在最大美德上卓越的本性，既不应是恶意的，也应当是善意的。因此，所有哲学家在义怒问题上意见一致，却在慈爱问题上分歧。但为使我的论述按顺序进入预定主题，我必须作出并遵循这类区分，因为义怒与慈爱是不同的，且彼此对立。要么将义怒归于天主，而将慈爱从祂夺去；要么两者都从祂夺去；要么将义怒除去，而将慈爱归于祂；要么两者都不除去。事情的本质除此之外再无其他可能；因此，所寻求的真理必然在这些中的某一个中被找到。让我们分别考虑它们，以便理性和次序将我们引向真理的隐匿之处。\par

% structure: p0011 source=h2 target=section reason=relative-heading-rank; base=h2

\section{第三章——论人类事务中的善与恶，及其作者}

首先，从未有人这样论及天主，说祂只易发怒，而不受慈爱的影响。因为赋予天主这样一种能力是不合适的：祂可能伤害并作恶，却不能施益与行善。那么，如果天主只是恶的作者，为人提供的是什么手段、什么得救的希望呢？因为倘若如此，那可敬畏的威严便将沦落，不是出于审判者的权能（审判者被允许保护和释放），而是出于施刑者和刽子手的职分。然而，既然我们看到人类事务中不仅有恶，也有善，那么显然，如果天主是恶的作者，就必须有另一位与天主相反、赐给我们善事者。如果存在这样一位，祂应被称为什么名字？为什么那伤害我们的，比那施惠于我们的更为我们所知？但如果这只能是出于天主，那么认为神圣的能力（没有比这更大或更好的）能够伤害却不能施益，就是荒谬而虚妄的；因此从未有人敢于断言此事，因为它既不合理，也绝不可信。既然这一点已被认同，让我们继续前往别处寻求真理。\par

% structure: p0013 source=h2 target=section reason=relative-heading-rank; base=h2

\section{第四章——论天主及其情感，以及对\Person{厄璧鸠鲁}的谴责}

接下来是关于\Person{厄璧鸠鲁}学派的观点：既然天主没有义怒，那么也确实没有慈爱。因为当\Person{厄璧鸠鲁}认为，伤害和施加损害（这大多源自义怒的情感）与天主不相称时，他也夺去了祂的善行，因为他看到，如果天主义怒，祂也必然有慈爱。因此，为了避免将恶习归于祂，他剥夺了祂的美德。由此，他说，祂是幸福且不朽坏的，因为祂什么也不关心，既不自寻烦恼，也不给别人带来烦恼。因此，祂不是天主，如果祂既不感动（这是活物的特性），也不做任何对人不可能的事（这是天主的特性），如果祂根本没有意志、没有行动，简言之，没有配得天主的管理。有什么比管理世界（尤其是人类，地上万物都受其管辖）更伟大、更相称的管理能归于天主呢？\par

那么，如果天主总是无所作为、安息而不动，祂有什么幸福可言呢？如果祂对祈祷者充耳不闻，对敬拜者视而不见，那又如何？有什么比眷顾更配得天主、更适宜祂呢？但若祂什么也不关心，什么也不预见，祂就完全失去了神性。那从天主那里夺去一切能力和一切实体的人，除了说根本没有天主之外，还说了什么？简言之，\Person{马尔库斯·图利乌斯}记载，\Person{波西多纽斯}曾说，\Person{厄璧鸠鲁}明白并没有神，但他所说关于神的话是为了避嫌；因此他在口头上保留了神，但实际上取消了神，因为他没有给他们任何运动、任何职责。但如果真是这样，还有谁比他更诡诈呢？这应当与智者和有分量的人的性格相去甚远。但如果他明白一事却说出另一事，除了称他为欺骗者、两面舌、邪恶者，且更是愚昧者之外，还能称他为什么？但\Person{厄璧鸠鲁}并非如此狡猾，以至于出于欺骗的欲望说那些话，因为他将这些事也写下来留作永久的纪念；他是由于对真理的无知而犯错。因为从一开始，他就被一个单一意见的表面真实性所引导，必然陷入那些随之而来的事。第一个意见是，义怒与天主的性格不相容。而这在他看来是真实且不可动摇的，他便无法拒绝其后果；因为一种情感被除去，必然性本身迫使他也从天主那里除去其他情感。因此，不受义怒支配的人显然也不受慈爱（与义怒相反的情感）的影响。现在，如果祂既无义怒也无慈爱，那么显然祂也没有恐惧、喜乐、忧伤或怜悯。因为所有情感都有同一个系统、同一个动向，而这不可能存在于天主。但如果天主没有任何情感，因为凡有情感的皆是软弱的，那么结论就是，祂也没有任何关心或眷顾。\par

智者的辩论至此为止：他对随后的事情保持缄默，即因为祂既无眷顾也无上智安排，所以祂内没有反思，也没有任何知觉，由此导致祂根本不存在。因此，当他逐渐下降时，他停留在最后一步，因为他现在看到了悬崖。但保持沉默并隐藏危险有何益处呢？必要性迫使他违背自己的意愿坠落。因为他说了他本无意说的话，因为他如此安排他的论证，以至于他必然到达他想要避免的那一点。因此，你看，当义怒从天主身上被移除和取走时，他到达了哪一点。简而言之，要么没有人相信这一点，要么只有极少数人，而且是那些有罪和邪恶的人，他们希望自己的罪不受惩罚。但如果这也被证明是假的，即天主既无义怒也无仁慈，那么让我们进入排在第三位的问题。\par

% structure: p0017 source=h2 target=section reason=relative-heading-rank; base=h2

\section{第五章　斯多葛学派关于天主的观点：论祂的义怒与仁慈。}

斯多葛学派和其他一些人被认为对天主本性持有更好的看法，他们说天主有仁慈，但没有义怒。一个非常悦耳且流行的说法是，天主不会如此心胸狭窄，以致认为有人伤害了祂，因为祂不可能被伤害；因此，那宁静而神圣的威严会被激动、烦扰和狂怒，这是人性的软弱所致。因为他们说，义怒是心灵的骚动和扰动，这与天主不相称。因为当它落到任何人的心灵上时，如一阵猛烈的风暴激起如此波浪，以致改变了心灵的状态：眼睛闪烁，面容颤栗，舌头结巴，牙齿打战，面容交替染上红晕和苍白。但如果义怒对一个有智慧和权威的人有失体统，那么如此可憎的变化对天主来说岂不更不相称！而且，如果一个人有权威和能力时，因义怒造成广泛伤害，流血、颠覆城市、摧毁社区、使省份沦为荒凉，那么，既然天主对全人类乃至宇宙本身拥有权能，如果祂发怒，岂不更应相信祂会毁灭万物！\par

因此他们认为如此巨大而有害的邪恶应当远离祂。如果义怒和激动远离祂，因为这是丑陋且有害的，祂不伤害任何人，那么他们认为别无他物，只余祂是温和、宁静、慈惠、慈善、保护者。因为只有这样，祂才能被称为万有的共同父亲，至善至大，这正是祂神圣而属天的本性所要求的。因为如果在人中间，行善似乎比伤害更值得称赞，使人活比杀人更值得称赞，拯救比毁灭更值得称赞，而无辜不被不公正地列为美德之一——而做这些事的人被爱戴、尊重、尊崇，并以一切祝福和誓言颂扬——简言之，因其功德和恩惠被认为最像天主；那么，天主自己，祂在神圣而完美的美德上卓越，且远离一切尘世污秽，岂不更应该以神圣而属天的恩惠安抚全人类！这些话说得动听且流行，引诱许多人相信；但持有这些观点的人确实更接近真理，但他们部分地失败了，没有充分考虑事情的本性。因为如果天主不向不虔敬和不义的人发怒，那么显然祂也不爱虔敬和正义的人。因此，那些同时剥夺义怒和仁慈的人的错误更为一贯。因为在相反之事上，必然要么牵动两面，要么都不牵动。因此，爱善的人也恨恶人，而不恨恶人的人也不爱善人；因为爱善源自恨恶，而恨恶源于爱善。没有不爱生命而恨死亡的人，也没有渴望光明而不避开黑暗的人。这些事情在本性上如此关联，以致一个不能没有另一个而存在。\par

若某家主有一好仆与一坏仆，显然他并非同样憎恨二者，亦非赐予二者同样恩惠与尊荣；因为若他这样做，他既是不义又是愚蠢。但他以友善之言对待那好仆，尊荣他，派他管理自己的家和家人，以及一切事务；却以斥责、鞭打、赤身、饥饿、口渴、锁链惩罚那坏仆：使后者作为他人戒除犯罪的榜样，前者则作为劝勉他人；好让恐惧约束一些人，尊荣激励另一些人。因此，那爱者也恨，那恨者也爱；因为有些人是应当爱的，有些人是应当恨的。正如爱者将善事赐予他所爱的，同样恨者将恶事加于他所恨的；这个论据，因为是真实的，绝无可反驳之处。因此，那些人的观点是虚妄而错误的：他们一方面将\OriginalTerm{仁慈}{kindness}归给天主，却否定\OriginalTerm{愤怒}{anger}；这与那些将两者都否定的人同样错误。但后者，如我们所示，并未全然犯错，而是保留了二者中更好的一面；而前者，因推理方法精确而受引导，落入最大的错误，因为他们采纳了全然虚假的前提。他们本不应如此推理：因为天主不受愤怒支配，所以祂也不受仁慈感动；而应如此：因为天主受仁慈感动，所以祂也受愤怒支配。因为若天主不受愤怒支配是确定而无疑的，那么另一论点必随之而来。但既然天主是否发怒更为可疑，而祂是仁慈的几乎完全显然，以不确定之事推翻确定之事是荒谬的，因为以确定之事证实不确定之事更为容易。\par

% structure: p0021 source=h2 target=section reason=relative-heading-rank; base=h2

\section{第六章——天主是愤怒的}

这些是哲学家们关于天主的见解。但若我们已发现所说的这些是虚假的，则剩下最后一条出路，唯有在此才能找到真理，它从未被哲学家采纳，也从未得到辩护：即必须得出结论：天主是愤怒的，因为祂被仁慈所感动。我们应持守并主张这一见解；因为这是全部虔敬与宗教所依赖的总纲与枢纽：若天主对敬拜者无所赐予，便无尊荣可归；若祂不向不敬拜者发怒，便无畏惧。\par

% structure: p0023 source=h2 target=section reason=relative-heading-rank; base=h2

\section{第七章——论人、禽兽与宗教}

虽然哲学家们常因对真理的无知而偏离理性，陷入无法解决的错误（因为那些不知道路又不承认自己不知的旅人，往往徘徊不定，羞于向遇见的人问路），但从未有哲学家宣称人与禽兽之间毫无区别。凡是希望显得有智慧的人，都未曾将理性动物降低到哑巴无理性的水平；而某些无知之人却如此行，他们如同禽兽本身，为放纵口腹与享乐，声称人与一切活物的受造原则相同，人说这样的话是亵渎的。因为谁如此无学问，谁如此缺乏理解，竟不知道人内有某种神圣之物？我尚未论及灵魂与理智的卓越，人由此与天主有显明的相似。难道身体姿态与面容本身不表明我们并非与哑巴造物同等吗？它们的本性俯向地面与食料，与天空无涉，因为它们不看天空。但人以直立姿态、高扬的面容朝向宇宙的默观，将自己的容貌与天主相比，理性认出理性。\par

因此，就如西塞罗所言，除了人，没有一种动物认识天主。因为惟独人具有\OriginalTerm{智慧}{wisdom}，所以他独能理解宗教；这是人与哑巴禽兽之间主要的或惟一的区别。因为其他看似人类特有的事物，即使在哑巴动物中不存在，却也可能看似类似。语言是人类特有的；然而在它们之中也有某种类似语言的东西。因为它们以声音互相区别；当它们发怒时，发出类似争吵的声音；当相隔一段时间后相见时，它们用声音表示庆贺。对我们而言，它们的声音确实粗野，正如我们的声音或许对它们而言也是如此；但对它们自己，因为它们互相理解，这就是话语。简言之，在每个情感中，它们发出不同的声音表现其心境。笑声也是人类特有的；然而我们在其他动物身上也看到某些喜乐的迹象，当它们以热情的姿态玩耍、垂耳、缩唇、展额、放松眼睛以示嬉戏时。有什么比\OriginalTerm{理性}{reason}和预见未来更为人所特有？但有些动物从巢穴开辟数个不同方向的出口，以便在危险来临时，被困的它们有逃脱之路；但若非它们具有智识与思考，是不会这样做的。另一些动物则预防未来，如\par

\Quote{蚂蚁抢劫一大堆谷粒，念念不忘冬季，将其储存在住所中；}\par

再者——\par

\Quote{又如蜜蜂，它们独知乡土与固定住所；且念及将临的冬季，夏季即劳作，将所得积聚为公共储备。}\par

若我欲追溯各禽兽族群通常所做的最似人工技艺之事，这将是一项冗长的任务。然而，若在所有这些通常归于人的事物中，连哑巴动物也发现有某种相似之处，那么显然，宗教是唯一在哑巴动物中找不到丝毫痕迹或迹象的事物。因为正义专属宗教，而其他动物无法达到。惟独人能统治；其他动物则服从于他。但天主的崇拜归于正义；凡不接纳此者，远离人性，将以人的形态过着禽兽的生活。然而，由于我们与其他动物几乎仅在此方面不同：即唯独我们感知天主的威能与权能，而其他动物没有对天主的认识，那么，在此方面，哑巴动物当然不可能更有智慧，人性也不可能愚钝，因为所有活物及整个自然体系因人的智慧而服从于人。因此，若理性、若人在此方面的力量超越并胜过其余活物，唯独他能认识天主，那么显然，宗教绝不能被推翻。\par

% structure: p0030 source=h2 target=section reason=relative-heading-rank; base=h2

\section{第八章——论宗教。}

但如果我们相信\Person{伊壁鸠鲁}这样说，宗教就被推翻了：——\par

\Quote{诸神的本性必然永远自享不朽与至高宁静，远离我们的关切，完全超脱；因为免于一切痛苦，免于一切危险，自足于自身资源，不需求我们任何东西，既不被恩惠所获，也不被愤怒所动。}\par

现在，当他说这些话时，他是否认为应向天主献上任何敬拜，抑或他完全推翻了宗教？因为若天主不赐予任何人任何善，若祂不以恩惠回报祂崇拜者的顺从，那么，还有什么比建造圣殿、奉献祭品、呈献礼物、减少我们的财产而一无所获更无意义、更愚蠢的呢？但（有人说）一位卓越的本性理应受到尊敬。对一位不关心我们、忘恩负义的存在，能给予什么尊敬？我们如何能与一位与我们毫无共同之处者有任何关联？\Person{西塞罗}说：\Quote{愿天主远离，}\Quote{若祂是如此，不受任何恩惠和人类情感的影响。因为为何我要说\InnerQuote{愿祂垂顾}？因为祂不能对任何人垂顾。}还有什么比这更蔑视天主的话呢？愿天主远离，他说，即让祂离开、退去，因为祂不能有益于任何人。但若天主不劳苦，也不给别人带来劳苦，那么，为何我们不在每次能避人耳目、欺骗公共法律时犯罪呢？只要我们获得有利的逃脱机会，让我们利用时机：让我们夺取他人的财产，或不流血，甚至若唯有法律值得敬畏，就流血也在所不惜。\par

\Person{伊壁鸠鲁}抱有这些想法时，他完全摧毁了宗教；一旦宗教被除去，生活的混乱与动荡将随之而来。但若宗教不能被除去而不破坏我们对智慧（藉此我们与禽兽分离）以及正义（藉此公共生活可更安全）的把握，那么，宗教本身如何能在没有恐惧的情况下被维持或守护呢？因为不被敬畏的事物被轻视，被轻视的显然不受尊敬。由此，宗教、威严与尊荣与恐惧共存；但若无人发怒，则无恐惧。因此，无论你从天主身上除去仁慈，还是愤怒，或二者兼有，宗教必然被除去；没有宗教，人的生活充满愚昧、邪恶与不法。因为良心大大约束人，如果我们相信我们在天主眼前生活；如果我们想象不仅我们所做的行为被从天俯瞰，而且我们的思想和言语也被天主听到。但相信这一点是有益的，如有些人所想象，并非为了真理，而是为了效用，因为法律无法惩罚良心，除非有来自上空的恐惧悬在头上以遏制罪行。因此，宗教全然虚假，并没有神性；所有一切都是巧手之人编造出来的，为使人生活得更正直、更无罪。这是一个重大问题，且与我们提出的主题无关；但因它必然出现，所以无论如何应当稍加处理。\par

% structure: p0035 source=h2 target=section reason=relative-heading-rank; base=h2

\section{第九章——论天主的眷顾，及反对它的意见。}

当昔日的哲学家们就眷顾达成一致意见，且毫不怀疑世界由天主和理性安置并受理性治理时，\Person{普罗泰戈拉}在\Person{苏格拉底}时代是第一个说他不清楚是否有任何神性的人。他的这一论辩被认为如此不敬、如此违背真理与宗教，以致\Person{雅典人}将他逐出领地，并在公众集会上焚烧了他载有这些陈述的书籍。但无需讨论他的意见，因为他没有说出任何确定的东西。此后，\Person{苏格拉底}及其门徒\Person{柏拉图}，以及从柏拉图学派如溪流般流向不同方向的那些人，即\Person{斯多葛学派}与\Person{逍遥学派}，与他们的前辈意见相同。\par

后来，伊壁鸠鲁说确实有一位天主，因为世界中必须存在某种卓越超群、非凡而蒙福的存在；然而并没有眷顾，因此世界本身并非按任何计划、技艺或工艺而被安排，而是由某种微小而不可分的种子构成了宇宙。但我看不出还有什么比这更违背真理的说法了。因为如果有一位天主，祂作为天主显然是眷顾的；并且除非祂保存过去、知晓现在、预知未来，否则无法以任何其他方式将神性归于祂。因此，在取消眷顾的同时，他也否认了天主的存在。但是，当他公开承认天主的存在时，他也同时承认了祂的眷顾，因为二者不能脱离对方而存在或被理解。但在后来的时期，哲学已失去活力，有一位米洛斯的狄亚戈拉斯，他完全否认天主的存在，并因这一见解被称为无神论者；还有昔兰尼的狄奥多若：这两人，由于无法发现任何新东西——一切都已被言说和发现——宁愿甚至违背真理，去否认所有先前的哲学家都毫无歧义地同意的观点。这些人攻击了眷顾，而这眷顾已被如此多的才智之士历经如此多的时代所主张和捍卫。那么如何呢？我们是要以理性，还是以杰出人物的权威，或者兼用二者来驳斥那些琐碎而怠惰的哲学家呢？但我们必须赶快前进，免得我们的言辞离题太远。\par

% structure: p0038 source=h2 target=section reason=relative-heading-rank; base=h2

\section{第十章——论世界的起源、事物的本性以及天主的眷顾}

那些不承认世界是由天主眷顾造成的人，要么说世界是由\OriginalTerm{第一原理}{first principles}偶然聚合而成，要么说世界是突然由自然产生的，但他们像斯特拉托那样认为，自然本身具有产生和减损的能力，但它既无感觉也无形状，这样我们就可以理解为，万物是自发地产生的，没有任何工匠或创造者。每一种观点都是虚妄且不可能的。但这发生在那些不认识真理的人身上，以至于他们杜撰出任何东西，也不愿认识到事物本性所要求的东西。首先，关于那些微小的种子，他们说是由这些种子的聚合形成了整个宇宙，我问这些种子在哪里或来自哪里。有谁曾见过它们？有谁曾感知过它们？有谁曾听说过它们？难道只有留基伯有眼睛吗？难道只有他有理智吗？而他肯定是所有人中唯一盲目和愚钝的，因为他所说的话，没有一个病人能在谵妄中说出，也没有一个睡着的人能在梦中说出。\par

古代哲学家们认为万物由四种元素构成。他不愿承认这一点，以免显得步人后尘；但他认为元素本身另有第一原理，这些原理既不能看见，也不能触摸，也不能被身体的任何部分感知。他说它们如此微小，以至于没有刀刃锋利到能切割和分割它们。因此他给它们取名为\OriginalTerm{原子}{atoms}。但他想到，如果它们全都具有同一种本性，它们就不可能构成我们在世界中所见的如此多样的不同物体。因此，他说有光滑的和粗糙的，有圆形的和多角的，还有带钩的。与其为如此可怜而空洞的用途而拥有舌头，保持沉默不是好得多吗？的确，我担心，凡认为这些事值得驳斥的人，似乎也同样是狂乱。然而，让我们像对一位说了某些话的人那样来回应。如果它们是柔软而圆形的，显然它们不能相互勾连以构成某个物体；就像有人想把小米粒粘合在一起，谷粒本身的柔软性不允许它们聚集成团。如果它们是粗糙的、多角的、带钩的，以便能够粘合，那么它们就是可分的，能够被切割；因为钩子和角必定突出，因此有可能被切掉。\par

因此，能够被切掉和撕裂的东西，也就能够被看见和握住。\Quote{这些，}他说，\Quote{以不停的运动在虚空中飞舞，到处飘动，正如当太阳通过窗户引入光线时，我们看见尘土微粒。从这些产生了树木、花草和地上的一切果实；从这些产生了动物、水、火和万物，它们又再次分解为相同的元素。}只要探讨的是微小事物，这一点还可忍受。就连世界本身也是由这些组成的。他已经达到了完全疯狂的地步：似乎不可能再说什么，但他还是找到了补充。\Quote{既然万物，}他说，\Quote{是无限的，而没有任何东西可以是空的，那么必然存在着无数个世界。}原子的力量竟如此巨大，以致从如此微小的元素中聚集起如此不可计量的质量？首先我追问：那些种子的本性或起源是什么？因为如果万物都来自它们，那么它们本身又来自哪里？什么本性为制造无数世界提供了如此丰富的大量物质？但让我们容许他逍遥地胡言乱语关于世界的事情；让我们谈论我们所在的这个世界，并谈论我们所看见的这个世界。他说万物都是由不可分割的微小物体构成的。\par

如果事情果真是这样，就没有任何物体需要同类的种子。鸟类将无卵而生，或卵而无出；其他生物亦然，无需交配；树木及大地的产物也将没有我们每日亲手处理并播种的种子。为何谷粒生出谷田，而谷田又生出谷粒？简言之，如果原子的相遇与聚合能产生一切，那么一切都会在空气中生长，因为原子在虚空中飘荡。为何香草、树木或谷粒，没有土壤、没有根、没有水分、没有种子，就不能生长或繁衍？由此可知，没有一样东西是由原子构成的，因为每样事物都有其独特而固定的本性、自己的种子、以及从太初就赋予的律法。最后，\Person{卢克莱修}——仿佛忘记了他所主张的原子的学说——为了反驳那些认为万物从无中产生的人，竟采用了这些可能反过来反对他自己的论据。因为他这样说：\par

\Quote{如果万物从无中来，任何种类都可能生于任何事物；没有东西需要种子。}\par

随后他又说：\par

\Quote{我们必须承认，无中不能生有，因为事物在各自诞生并进入广袤的空气中之前，都需要种子。}\par

谁会想到他说这些话时是有头脑的，却没有看出它们自相矛盾？因为从每样事物都有确定的种子这一点来看，没有东西是由原子造成的——除非我们偶然相信火和水的本性也来自原子。为何我要说，如果最坚硬的物质猛烈撞击，就会迸出火花？难道原子藏在钢中，或燧石里？谁将它们关在里面？为何它们不自由迸出？火的种子如何能留在最寒冷的物质中？\par

我暂且不提燧石与钢。如果你在阳光下拿着一个装满水的晶球，即使在最严寒的天气里，水反射的光也会点燃火。那么我们必须相信水中含有火吗？然而即使在夏天，太阳也不能从水中生火。如果你向蜡呼气，或一股轻气触及任何东西——无论是坚硬的大理石表面还是金属板——水都会以极小的水滴逐渐凝结。同样，从大地或海洋的蒸发形成雾气，它或扩散开来，湿润所覆盖的一切；或聚集起来，被风带到高山之巅，压成云层，降下大雨。那么，我们说流体是在哪里产生的？是在水汽中？在蒸发中？在风中？但任何既不能被触摸也不能被看见的东西中，什么也不能形成。为何我要谈论动物？在它们的身体中，我们看不到任何没有计划、没有安排、没有用途、没有美感而形成的东西，以至于所有部分和器官的巧妙而周密的配置都排除了偶然与碰巧的可能性。但让我们假设肢、骨、神经和血液可以由原子构成。那么感觉、反思、记忆、心智、天赋的才能——它们能用什么种子聚合起来？他说，来自最微小的种子。因此还有更大的种子。那么它们又如何是不可分的呢？\par

其次，如果不可见的事物由不可见的种子形成，那么可见的事物就应由可见的种子形成。那么为何没有人看见它们？但无论一个人注视人身上的不可见部分，还是可触摸且可见的部分，谁看不出这两部分都是按照设计而存在的？那么，无目的地相遇的物体如何能产生任何有理性的东西？因为我们看见，在整个世界中没有任何事物自身不具有非常伟大奇妙的计划。既然这超越人的感官与能力，除了归因于天主的眷顾，还有什么是更正确的呢？如果一尊肖似人的雕像是由设计与技艺所造，难道我们设想人本身是由随意聚集的碎片构成的吗？而且被造之物与真理有何相似呢？因为最伟大、最卓越的技艺也只不过能模仿身体的粗略轮廓和外表线条。人所能赋予其作品的，有何种运动或感觉？我不提及视觉、听觉、嗅觉的运作，以及其他器官——无论是显露在外的还是隐藏不见的——的奇妙功用。哪个工匠能造出人的心脏、声音或智慧本身？因此，任何心智健全的人会认为，靠理性与判断人不能完成的事，却可以由到处相互粘附的原子的相遇来实现吗？你看他们陷入了何等愚蠢的胡言乱语，因为他们不愿将万物的创造与眷顾归于天主。\par

但让我们允许他们，世俗之物是由原子构成的：那么天上的事物也是吗？他们说诸神是无染污的、永恒的、有福的；他们唯独豁免了诸神，使得他们似乎不是由原子的集合而成的。因为如果诸神也是由这些构成的，他们就会容易消散，最终种子分解并回归其本性。因此，如果有些东西原子无法产生，我们为何不能以同样的方式判断其他事物？但我问，为什么诸神在那些最初的元素产生世界之前，没有为自己建造一个居所？显然，除非原子聚集并形成了天，否则诸神仍将悬浮于虚空之中。那么，凭着何种谋略，凭着何种计划，原子从混乱的团块中聚集起来，使得从某些原子中，下面的地形成球体，而天在上面展开，装饰着如此众多的星座，以至于无法想象更精美的东西？因此，看见如此伟大事物的人，怎能想象它们是在没有任何设计、没有任何眷顾、没有任何神圣智慧的情况下形成的，而是由细微的原子产生出如此伟大的奇妙事物？这岂不像是奇闻，竟然存在能说出这些话的人，或存在相信这些话的人——如\Person{德谟克利特}，他是他的听众，或\Person{伊壁鸠鲁}，一切愚昧都从\Person{留基伯}的泉源流向他？但是，正如其他人所说，世界是由自然——没有知觉和形状——所造的。但这更加荒谬。如果自然造了世界，它必定是凭着判断和智慧造的；因为造物者必须有造物的倾向或知识。如果自然没有知觉和形状，那么有知觉和形状的事物怎能由它造出？除非有人偶然认为，如此精巧的动物构造，竟能由无知觉之物形成并赋予生命，或者那个为生物之用而如此预见地准备的天穹，竟在无建造者、无工匠的情况下突然出现。\par

\Quote{如果有什么东西，}\Person{克里斯普}说，\Quote{能成就那些有理性的人所不能成就的事，那东西必定比人更大、更强、更有智慧。}但人不能造天上的事物；因此，产生或已经产生这些事物的东西，在技艺、设计、技能和能力上都超越人。那么，除了天主，还能是谁呢？但是，他们视为万物之母的自然，如果没有理智，将一事无成，毫无设计；因为没有思想的地方就没有运动或效力。但如果它运用谋略来开始任何事情，运用理性来安排，运用技艺来完成，运用能量来实现，并运用能力来统治和控制，那么为何称它为自然而不称它为天主？或者，如果原子的集合或无思想的自然造出了我们所看见的事物，我问它为何能造天，却不能造城市或房屋？为何它造了大理石山，却没有造柱子和雕像？既然原子尝试所有位置，难道它们不该聚集起来做成这些东西吗？至于没有思想的自然，它忘记做这些事并不奇怪。那么，情况如何？显然，天主在开始这世界的工作时——没有比它更有秩序、更适应用途、更美丽、更伟大的了——祂自己造了人不能造的东西；其中也有人自己，祂给了人一份祂自己的智慧，并赋予他理性，以配合地上的脆弱所能接受的，以便他能为自己制造必需品。\par

但是，如果在这世界的共同体——姑且这么说——中，没有统治的眷顾，没有治理的天主，那么在这事物的本性中便没有任何知觉。那么，从何处相信人的心灵及其技艺和智慧有其起源呢？因为如果人的身体是由土造成的（人因此得名），那么灵魂——有智慧，是身体的主宰，四肢像服从君王和统帅一样服从它，既不能看见也不能理解——不可能来自除智慧本性之外的其他东西。但正如心灵和灵魂管理每一个身体，同样天主也管理世界。因为不太可能的是，较小和卑微的东西拥有统治权，而较大和最高的东西却没有。简言之，\Person{马尔库斯·西塞罗}在他的\WorkTitle{图斯库卢姆辩论集}和\WorkTitle{安慰}中说：\Quote{灵魂的起源在地上找不到。因为他说，灵魂中没有混合和复合的东西，也没有看起来由土产生和构成的东西；没有湿润或空气或火性质的东西。因为这些本性中没有记忆、心灵和反思的力量，这些力量既保留过去又预见未来，并能把握现在；只有这些东西是神圣的。因为永远找不到任何源泉，除了来自天主，它们能够来到人身上。}因此，除了两三个虚妄的诽谤者外，人们一致认为世界由眷顾治理，也如此被造，没有人敢将\Person{狄亚戈拉斯}和\Person{狄奥多鲁斯}的意见，或\Person{留基伯}的空洞虚构，或\Person{德谟克利特}和\Person{伊壁鸠鲁}的轻率，置于那七位被称为智者的古代人物，或\Person{毕达哥拉斯}、\Person{苏格拉底}、\Person{柏拉图}以及其他认为存在眷顾的哲学家的权威之上；因此，那一意见也是错误的，即他们认为宗教是由智者为了恐吓和恐惧而设立的，以便无知的人能避免犯罪。\par

但若此事属实，则我们受古代智者愚弄。然而，若他们为欺骗我们而发明宗教，更欺骗全人类，则他们并非智者，因为谎言与智者的品格不符。但且承认他们是智者；他们竟能欺骗不仅无知者，更欺骗\Person{柏拉图}、\Person{苏格拉底}，并轻易迷惑\Person{毕达哥拉斯}、\Person{芝诺}、\Person{亚里士多德}这些最大派别的首领，这是何等巨大的欺骗成就？故有天主眷顾，正如我所提及的那些人所觉察的，藉其运行与能力，我们所见一切事物皆被造并受治理。如此庞大的体系、如此安排与秩序以保存固定的次序与时间，若无一有眷顾的工匠，则起初无法产生；若无一有力的居住者，则不能存在如此多世纪；若无一熟练智慧的管理者，则无法永久治理；理性本身也宣告此事。凡有理性之物，必出于理性。而理性属于智慧而有理智的本性；智慧而有理智的本性无非是天主。世界既有理性，由之治理并凝聚，故为天主所造。若天主是世界的创造者与统治者，则宗教正确而真实地建立；因为尊荣与敬拜应当归于万物的创造者与共同父亲。\par

% structure: p0053 source=h2 target=section reason=relative-heading-rank; base=h2

\section{第11章——论天主，独一的天主，世界由其眷顾治理而存在。}

既然关于眷顾已达成一致，接着要表明的是，应相信眷顾属于多位，还是只属一位。我认为，我们在\WorkTitle{训导}中已充分教导，不能有众多神祇；因为若神圣的德能与权力分散于多位，则必被削弱。但被削弱的即是可朽的；若祂非可朽，则既不能被削弱也不能被分割。故只有一位天主，其圆满德能既不能被削弱也不能被增加。但若有众多神祇，当他们各自拥有一部分权力与权威时，总和本身即减少；他们不能各自拥有与别人共享的全体：每位所缺的正是别人所拥有的。故此，这世界不能有众多统治者，一屋不能有众多主人，一船不能有众多舵手，一群或一队不能有众多首领，一蜂群不能有众多蜂王。天上亦不能有众多太阳，如同一个身体内不能有多个灵魂；整个自然完全一致于统一。但若世界\par

\Quote{被一个灵魂滋养，\newline{}一股灵性的火焰，\newline{}在肢体的每一部分发光，\newline{}并搅动巨大的整体。}\par

由诗人的见证可知，有一位天主居住世界，因整个身体只能被一个心智居住治理。故所有神圣权力必须集中于一位，由祂的意志与命令统治万物；因此祂如此伟大，人不能用言语描述，也不能凭感官测度。那么，关于众多神祇的意见或信念从何而来？无疑，那些被尊为神的都是人，且是最早最伟大的君王；但谁不知道，他们死后被授予神圣尊荣，或因他们有益于人类的德行，或因他们的恩惠与发明使人生活增辉而获得不朽记忆？不仅男人，也有女人。对此，希腊最早的作家（他们称为神学家），以及追随模仿希腊人的罗马作家都教导；其中最著名的是\Person{欧赫墨洛斯}与我们的\Person{恩尼乌斯}，他们指出了所有神祇的生日、婚姻、后代、政绩、功业、死亡与坟墓。\Person{图略斯}追随他们，在其第三卷\WorkTitle{论神性}中摧毁了公共宗教；但无论是他还是其他人都未能引入真正的宗教，因为他对此无知。故他亲自作证，虚假的已显明，真理却隐藏。他说：\Quote{但愿上天，}他说，\Quote{我能轻易发现真相如同驳斥谬误！}他如此宣告，并非以学园派的掩饰，而是真诚地依其内心情感，因为真理不能从人类知觉中根除：人的预见所能达到的，他已达到，以便揭露虚假之事。因为一切虚构虚假之事，因无理性支持，容易摧毁。故只有一位天主，万物的泉源与本原，正如\Person{柏拉图}在\WorkTitle{蒂迈欧篇}中所感受并教导的，祂的伟大如此崇高，既不能为心智所理解，也不能为言语所表达。\par

\Person{赫尔墨斯}也作同样的见证，\Person{西塞罗}断言\Place{埃及人}将他列为众神之一。我说的是那位因其卓越和对多种技艺的知识而被称为\Person{特里斯墨吉斯托斯}的人；他不仅远早于\Person{柏拉图}，而且早于\Person{毕达哥拉斯}和那七位智者。在\Person{色诺芬}的著作中，\Person{苏格拉底}在谈话时说，不应探究天主的形象；\Person{柏拉图}在\WorkTitle{法律篇}中说：\Quote{天主是什么，不应成为探究的主题，因为既无法发现也无法讲述。}\Person{毕达哥拉斯}也承认只有一个天主，说有一个非物质的理智，它扩散并贯穿整个自然，赋予所有生物生命的知觉；但\Person{安提斯泰尼}在他的\WorkTitle{物理学}中说，只有一个自然的天主，尽管各民族和各城邦有他们自己人民的神。\Person{亚里士多德}及其逍遥学派，以及\Person{芝诺}及其斯多葛学派，所说的大致相同。确实，逐一追溯所有个人的观点会是一项冗长的工作，他们虽然使用不同的名称，但在治理世界的同一权能上是一致的。然而，尽管哲学家和诗人，以及简言之，那些崇拜诸神的人，常常承认至高天主，却从未有人探究过，从未有人讨论过，对祂的敬拜和尊崇之事；事实上，他们总是相信祂是仁慈和不朽坏的，认为祂不对任何人发怒，也不需要任何敬拜。因此，哪里没有恐惧，哪里就没有宗教。\par

% structure: p0058 source=h2 target=section reason=relative-heading-rank; base=h2

\section{第十二章——论宗教与敬畏天主。}

现在，既然我们已经回应了一些人的不敬和可恶的智慧，或者更确切地说，是愚蠢，让我们回到我们提出的主题。我们说过，如果剔除了宗教，智慧和正义都无法保持：智慧，因为对天主本性的理解（我们在这一点上不同于畜类）只存在于人身上；正义，因为除非那不可欺骗的天主约束我们的欲望，否则我们将邪恶和不敬地活着。因此，我们的行为被天主注视，不仅关系到共同生活的益处，甚至关系到真理；因为如果剔除了宗教和正义，失去理性，我们要么堕落到畜类的无知，要么堕落到野兽的野蛮，甚至更甚，因为野兽也饶恕自己的同类。还有什么比人更野蛮、更无情呢？如果对至高存在的恐惧被除去，他就能逃脱法律的注意或藐视法律的威力？因此，唯独对天主的恐惧护守着人与人之间的共处，生命本身也赖此得以维持、保护和管理。但是，如果人相信天主没有忿怒，这种恐惧就被除去了；因为天主在发生不义行为时会动怒和气愤，这不仅为共同利益所证明，而且理性本身和真理也说服我们。我们必须再次回到之前的话题，正如我们已教导世界是天主所造，我们也要教导为何而造。\par

% structure: p0060 source=h2 target=section reason=relative-heading-rank; base=h2

\section{第十三章——论世界的益处和功用及季节。}

如果有人思考整个世界的管理，他必定会明白斯多葛学派的观点是多么真实，他们说世界是为了我们而造的。因为世界所构成的一切，以及它自身所产生的一切，都适用于人的使用。因此，人用火来取暖、照明、软化食物和加工铁器；他使用泉水来饮用和沐浴；他使用河流来灌溉田地，并为国家划分边界；他使用大地来收获各种果实，用丘陵来种植葡萄园，用山地来获取树木和柴薪，用平原生产谷物；他使用海洋不仅为了贸易和从远方国家获取供应，也为了获取各种丰富的鱼类。但如果他使用这些与他最接近的元素，那么毫无疑问他也使用天空，因为甚至天上的事物也为了我们赖以生存的大地的肥沃而调节。太阳以其不断的运行和不等的间隔完成它的周年循环，或在升起时引出白天供劳作，或在落下时带来黑夜供休息；它时而向南远行，时而向北靠近，引起冬夏交替，以便冬季的湿气和霜冻使大地肥沃多产，夏季的炎热使草类的产物成熟硬化，或使潮湿地方的产物经过蒸煮和加热而成熟。月亮掌管夜晚的时间，以光明的盈亏调节她的月度周期，并以她的光辉驱散阴沉的黑暗照亮黑夜，使得在夏季炎热中的旅行、远征和工作可以没有辛苦和不便地进行；因为\par

\Quote{夜间割取轻巧的谷茬，夜间\newline{}干燥的草地更适于收割。}\par

其他天体，或在其升起或降落时，也以其固定的位置提供有利的时令。此外，它们还为船只提供引导，使船只不至于在无边的深渊中漫无目的地漂泊，因为舵手适当观察它们，就能抵达他所瞄准的岸边港口。云被风的气息所吸引，使播种的田地得以靠阵雨浇灌，使葡萄树硕果累累，树木果实丰盈。这些事是通过一年中变化更替而显现的，以便任何时候都不缺乏维持人类生命所需之物。但（有人说）同一大地也滋养其他生物，连哑巴牲畜也靠同样的产物喂养。难道天主也为了哑巴牲畜而劳作吗？绝不；因为它们没有理性。相反，我们明白，就连这些牲畜本身也是天主为供人使用而造的：一部分作为食物，一部分作为衣物，一部分协助人工作；由此可见，天主眷顾愿意以丰盛的事物和资源来供应和装饰人类的生活，因此祂以飞鸟充满天空，以鱼类充满海洋，以四足动物充满大地。然而，\OriginalTerm{学园派}{Academics}在与\OriginalTerm{斯多葛派}{Stoics}辩论时，惯常质问：如果天主为人类创造了万物，为什么在海中、在陆地上却有许多东西甚至与我们敌对、有害、伤害我们？而斯多葛派不顾真理，极其愚蠢地反驳这一点。因为他们说，在自然产物中，在动物之列，有许多东西的用途至今尚未被发现，但随着时间的推移，正如需要和用途已经发现了许多先前时代未知的事物一样，这些用途也会被揭示。那么，在老鼠、甲虫、蛇这些对人类麻烦且有害的东西中，能发现什么用途呢？是不是它们身上隐藏着某种药物？如果有，总有一天会被发现，即作为对抗邪恶的疗法，而他们抱怨说这完全是邪恶的。他们说，蝰蛇被烧成灰后，是治疗同一种动物咬伤的药物。与其需要从它自身提取的疗法，那它根本不存在岂不更好？\par

他们本可以更简洁、更真实地如此回答：当天主造了人，如同祂自己的肖像——这是祂做工的完成——祂唯独将智慧吹入人内，使人能把万物置于自己的权威和治理之下，并利用世间一切益处。然而，祂也将善与恶的事物摆在他面前，因为祂赐给了他智慧，而智慧的全部本性就在于辨别邪恶与良善的事物：因为除非人同时知道拒绝和避免邪恶的事物，否则他无法选择更好的事物并认识什么是善。两者相互关联，以致于若除去其一，另一也必然被除去。因此，既然善恶事物被摆在其前，智慧才终于履行其职责，为有用而渴望善，为安全而拒绝恶。因此，正如有无数的善事被赐予人，以便他能享用，同样也有恶事，以便他能防范。因为如果没有邪恶，没有危险——简言之，没有任何东西能伤害人——那么智慧的全部素材就被除去，对人而言便成为不必要的了。因为如果只有善事摆在眼前，哪里还需要思考、理解、知识和理性呢？因为无论他伸向何处，那都是合适且顺应自然的；这就好比有人想在尚未有辨味的婴儿面前摆上一顿极其精致的晚餐，显然每个婴儿都会渴望那被冲动、饥饿或甚至偶然所吸引之物；而无论他们取用什么，都将是有益和健康的。那么，他们永远保持现状、永远是婴儿、不谙世事，又有什么损害呢？但如果你在其中加入一些苦的、无用的，甚或有毒的东西，那么他们因无知于善恶而显然会被欺骗，除非将智慧加给他们，使他们能拒绝恶事、选择善事。\par

由此可见，我们因着罪恶而更需要智慧；若非这些事物被置于我们面前，我们就不会成为理性的受造物。但如果这个说法是真的——这是\OriginalTerm{斯多葛派}{Stoics}完全无法看见的——那么\Person{伊壁鸠鲁}的论证也归于无效。他说：\Quote{天主要么愿意除去罪恶却不能；要么能够却不愿意；要么既不愿意也不能；要么既愿意又能够。如果祂愿意却不能，祂便是无能的，这与天主的属性不符；如果祂能够却不愿意，祂便是嫉妒的，这同样与天主相悖；如果祂既不愿意也不能，祂便是既嫉妒又无能，因此就不是天主；如果祂既愿意又能够——这唯独与天主相称——那么罪恶从何而来？祂为何不除去它们？}我知道许多拥护天主眷顾的哲学家常被这个论证困扰，几乎违背自己的意愿而承认天主对万事不感兴趣，这正是\Person{伊壁鸠鲁}特别要达到的目的；但仔细考察后，我们轻易地化解了这个可怕的论证。因为天主能够做祂所愿的一切，天主内既无软弱也无嫉妒。因此祂能够除去罪恶；但祂不愿意这样做，然而祂并不因此是嫉妒的。因为祂不除去它们的原因，正如我已指出的，是祂同时赐予智慧；而智慧中的美善与喜乐，远胜于罪恶中的烦恼。因为智慧甚至使我们认识天主，并藉此认识获得永生，这是至高的善。因此，除非我们先认识罪恶，我们将无法认识善。但\Person{伊壁鸠鲁}没有看见这一点，也没有任何人看见：如果罪恶被除去，智慧也会同样被除去；人内便不再留下任何德行的痕迹，德行的本性正是在于忍受并克服罪恶的苦涩。因此，为了除去罪恶这一点微小的益处，我们竟会被剥夺一个极其伟大、真实且专属于我们的善。所以显然，万有都是为人而设，不论罪恶还是善好。\par

% structure: p0066 source=h2 target=section reason=relative-heading-rank; base=h2

\section{第十四章——天主为何造人。}

接下来我要指出的是，天主造人本身的目的。祂既为人的缘故而创造世界，同样也为祂自己的缘故而造人，使人如同神圣殿堂的司铎、祂工程与天上之物的观赏者。因为唯有人，由于他有理智、能推理，才能认识天主、赞叹祂的工程、觉察祂的能力与权能；为此，他蒙赐予判断、智力和谨慎。正因如此，他超越其他生灵，被造成直立的身姿与体态，仿佛被抬高以便仰望他的父。正因如此，唯独他领受了言语和作为思想通译的舌头，使他能宣扬祂主的威严。最后，为此万有都被置于他的管辖之下，而他自己则应在天主——万物的创造者与造物主——的管辖之下。因此，若天主设计人作为祂的敬拜者，并为此赐予他如此大的尊荣，使他统管万有；那么显然，人理应敬拜那赐予他如此大恩的祂，并爱那与我们共享神圣正义的人。因为一位天主的敬拜者不应被另一位天主的敬拜者伤害。由此可知，人是为了宗教和公义而被造的。关于此事，\Person{马库斯·图利乌斯}在其论法律的著作中作证，他这样说：\Quote{但在学者们争论的一切事物中，最重要的莫过于完全明白我们生来就是为了公义。}如果这是极为真确的，那么天主愿意所有的人都成为公义的，即爱天主和爱人如己；真诚地尊崇天主为父，并爱人为弟兄：因为全部公义涵盖在这两件事中。但不认识天主或损害他人的人，便生活得不公义，违反他的本性，从而扰乱了神圣的制度和律法。\par

% structure: p0068 source=h2 target=section reason=relative-heading-rank; base=h2

\section{第十五章——罪恶如何延伸至人。}

也许有人会问：罪恶如何延伸至人，或是何种扭曲将神圣制度的准则败坏得更糟，以致人虽生为公义，却行不义之事？我在前面已经解释过，天主同时将善与恶摆在他面前，祂爱善而恨与善相反的恶；但祂允许恶存在，为的是善也能彰显出来，因为我常教导说，二者缺一不可；简言之，世界本身由两种对立且相连的元素构成，即火与水，且若无黑暗，光也无法被造，因为没有低处就没有高处，没有日落就没有日出，没有寒冷就没有温暖，没有柔软就没有坚硬。同样，我们也由两种同等对立的实体构成——灵魂与身体：其中之一归属于天，因为它轻盈不可触摸；另一个归属于地，因为它可被把握：一个是坚固而永恒的，另一个是脆弱而必死的。因此，善依附于一个，恶依附于另一个：光明、生命和公义属于一个；黑暗、死亡和不义属于另一个。因而在人中产生了本性的败坏，以致有必要建立一条法律，藉以禁止恶习，并规定德行的义务。既然人的事务中有善与恶的事物，其本性我已阐明，那么天主必定对两方面都有所动：当祂看见正义之事被行时，祂便施恩；当祂察觉不义之事时，祂便发怒。\par

但厄壁鸠鲁反对我们，说：\Quote{若天主内有喜悦之情使其偏爱，有仇恨之情使其愤怒，则祂必然也有恐惧、倾向、欲望及属于人性软弱之其他情感。}这并非说，愤怒的人必恐惧，喜悦的人必悲伤；简言之，易于愤怒者较少胆怯，性情喜悦者较少受悲伤影响。何须谈论人性情感，我们的本性屈服于它们？让我们权衡天主的必然性；因为我不愿谈论本性，因为相信我们的天主从未受生。恐惧之情在人内有其对象，但在天主内则无。人，因易遭许多意外和危险，恐惧更大的暴力可能发生，袭击、掠夺、撕裂、击倒并毁灭他。但天主，既无匮乏、伤害、痛苦、死亡，绝不能恐惧，因为没有任何事物能对祂施以暴力。同样，欲望的理由和原因在人内是明显的。因为人受造脆弱而可朽，必须造出另一个不同性别，与之结合以产生后代延续种族。但此欲望在天主内无位置，因为脆弱和死亡远离祂；祂内也无女性可与结合以喜悦；祂无需承继，因祂永远活着。关于嫉妒和激情也可说同样的话，人有确定而明显的原因易受它们影响，但天主绝不如此。然而，事实上，偏爱、愤怒和怜悯在天主内有其实质，那最伟大无比的大能运用它们维护世界。\par

% structure: p0071 source=h2 target=section reason=relative-heading-rank; base=h2

\section{第十六章 论天主及其忿怒与情感}

有人会问这实质是什么。首先，当灾祸降临时，人在沮丧中大多会求助于天主：他们平息并恳求祂，相信祂能抵御伤害。因此祂有施怜悯的时机；因为祂并非如此冷酷和藐视人，以至于拒绝帮助困苦者。同样，许多深信正义悦乐天主的人，既敬拜那为万物之主与父的天主，又以不断的祈祷和重复的誓愿献上供物和牺牲，以赞美追随祂的名，努力以正义和善行获得祂的眷顾。因此有一个理由，天主可以且应当眷顾他们。因为若没有比仁慈更合宜于天主的事，也没有比忘恩负义更不适合祂品格的事，那么祂必须对那些卓越、度圣洁生活者的服务有所回报，以免犯下即使在人身也值得责备的忘恩负义之罪。但相反，另一些人胆大妄为、邪恶，以私欲玷污一切，以杀戮骚扰，行欺诈，掠夺，发假誓，不宽恕亲属和父母，藐视法律，甚至藐视天主本身。因此，愤怒在天主内有其合宜的时机。\par

因为当祂看到这样的事而不动心，起来报复恶人、消灭害人和有罪者以促进所有善人的利益，是不对的。因此，即使在愤怒本身之中，也包含着仁慈的显示。所以，那些既不承认天主会愤怒，又坚持祂显示仁慈（因为没有愤怒这确实不可能发生）之人的论证，以及那些认为天主内没有情感波动之人的论证，被发现是空洞和虚假的。由于有些情感（如欲望、恐惧、贪婪、悲伤和嫉妒）是天主所不具有的，他们便说天主完全没有情感。因为祂不具有这些，因为它们是邪恶的情感；但属于美德的情感——即对恶人的愤怒、对善人的眷顾、对受苦者的怜悯——因为配得上神圣能力，祂拥有自己正义而真实的情感。如果祂不拥有这些，人类生活将陷入混乱，事物状况将如此动荡，以致法律被藐视和压倒，唯独胆大妄为统治，以至于最终无人能安全，除非以力量超群者。这样，全地将如同被普遍的抢劫所破坏。但如今，既然恶人期待惩罚，善人盼望眷顾，受苦者寻求援助，美德便有空间，罪行更为稀少。但有人说，恶人往往更兴旺，善人更悲惨，义人被不义者肆意骚扰。这些事为何发生，我们将在以后考虑。与此同时，让我们解释关于愤怒：天主内是否有愤怒？祂是否完全不顾，对以不敬虔所做之事无动于衷？\par

% structure: p0074 source=h2 target=section reason=relative-heading-rank; base=h2

\section{第十七章 论天主的眷顾与忿怒}

\Person{伊壁鸠鲁}说，天主毫不关怀任何事；因此祂没有权能。因为凡有权能的，必然要处理事务。如果祂有权能，却不用它，那么有什么重大原因会使祂眼中——我且不说我们人类，甚至连宇宙本身——都显得可鄙呢？\newline{}为此，他说祂是纯洁而幸福的，因为祂常处于安息之中。那么，如此重大事务的管理权交给了谁，如果这些我们看见由最高智慧所治理的事物竟被天主所忽视？\newline{}又或，那位生活着、知觉着者，如何能安息？因为安息要么属于睡眠，要么属于死亡。但睡眠并非安息。当我们睡眠时，身体固然安息，灵魂却不安宁、骚动不息：它为自己构成可观看的影像，从而以种种幻象行使自身本性的运动能力，并将自己从虚假事物中唤回，直到肢体满足，并从安息中获得力量。\newline{}因此，永恒的安息唯独属于死亡。既然死亡不影响天主，那么天主就永不安息。\newline{}然而，天主的行动除治理世界之外，还能是什么？\newline{}但如果天主持续照料世界，那祂就必会关怀人类的生命，留意个人的行为，并热切期望他们成为智慧而良善的。这是天主的旨意，这是神圣的法律；遵守并持守这法律的人，为天主所爱。\newline{}祂必然会对那违背或轻视这永恒神圣法律的人动怒。\newline{}他说，如果天主伤害任何人，那祂就不是良善的。他们受不小的谬误所欺骗，他们用严厉与恶意的名称诋毁一切——无论是人类的还是天主的——责罚，认为那以惩罚报应伤人者的人应当被称为伤人者。\newline{}但倘若如此，那么惩罚罪犯的法律就是伤人的，对罪犯处以死刑的法官也是伤人的。\newline{}但如果那给违法者以应得之报的法律是正义的，如果那惩罚罪恶的法官被称为正直和良善的——因为惩治恶人就是保护善人的安全——那么，天主对抗恶人并非伤人；反之，那伤害无辜者，或容忍伤人者以伤害更多人的，才是伤人的。\par

我乐意问那些把天主描述为不动者的人：如果某人拥有财产、房屋、一群奴隶，而他的奴隶藐视主人的宽容，侵占一切，并自己享用他的财物，他的家人反倒尊崇他们，而主人却被众人藐视、侮辱、遗弃——那么，不报复侮辱，反而容许那些在他权下的人享用他的财物的人，能算是智者吗？\newline{}这样的宽容能在任何人身上找到吗？\newline{}如果这确实可以称作宽容，倒不如说是某种麻木不仁。\newline{}然而，忍受藐视是容易的。\newline{}但如果发生了\Person{西塞罗}所谈论的那些事呢？\newline{}我问：如果一位家长，他的孩子被一个奴隶杀死，妻子被杀，房屋被烧毁，他却不对那奴隶施加最严厉的惩罚，那么他会被视为仁慈和悲悯，还是无人性和残忍至极？\newline{}但如果宽恕这类行为与其说是仁慈，不如说是残忍，那么天主对不义之事不动怒就不是善性的一部分了。\newline{}因为世界仿佛是天的家，人类仿佛是祂的奴隶；如果祂的名成为他们的笑柄，那么放弃祂自己的尊荣，看到邪恶与不义之事发生而不愤怒——这本是祂那对罪愆不悦之本性所特有和自然的——这算哪种或多大程度的宽容？\newline{}因此，发怒是理性的一部分：因为这样，过错被消除，放纵被约束；这显然符合正义和智慧。\par

但斯多葛派没有看到对与错之间、正义与不义之怒之间的区别；因为他们找不到处理这问题的方法，就想彻底消除它。\newline{}而逍遥派说它不应被除掉，而应被调节；我们在\WorkTitle{制度}第六卷中已对此作了充分的回应。\newline{}如今，哲学家们对愤怒的本性无知，这从他们的定义可以清楚看出——\Person{塞内卡}在他关于愤怒主题所写的著作中列举了这些定义。他说：\Quote{愤怒是对伤害复仇的欲望。}\newline{}另一些人，如\Person{波西多尼乌斯}所说，将其描述为惩罚你认为曾对你施加不公伤害之人的欲望。\newline{}有些人这样定义：\Quote{愤怒是心灵的一种冲动，旨在伤害那已经施行了伤害或曾希望施行伤害的人。}\newline{}亚里士多德的定义与我们的定义相差不大；他说\Quote{愤怒是回报痛苦的欲望。}\newline{}这是不义的愤怒，就是我们之前提到的，甚至存在于哑巴动物中；但在人身上必须加以约束，以免他因暴怒而冲入某种极大的邪恶。\newline{}愤怒不能存在于天主内，因为祂不可能被伤害；但它存在于人内，因为人是脆弱的。因为伤害的施加燃起痛苦，而痛苦产生复仇的欲望。\newline{}那么，那针对罪人的正义之怒何在？因为这显然不是复仇的欲望，因为没有伤害在先。\newline{}我并非指那些违反法律的人；因为尽管法官对这些人生气无可指责，但我们还是假定他在判决罪犯受罚时应保持镇定，因为他是法律的执行者，而非自己心意或权力的执行者；这正是那些试图根除愤怒的人所希望的。\newline{}但我特别指那些在我们权力之下的人，如奴隶、孩童、妻子和学生；因为当我们看到这些人犯错时，我们就被激起去约束他们。\par

因为那正义而良善者必定对邪恶之事感到不悦，而对邪恶感到不悦者在看到邪恶被施行时必会激动。因此我们起身报复，不是因为受到了伤害，而是为了保存纪律、纠正道德、抑制放纵。这是正当的愤怒；正如它在人身上对于纠正邪恶是必要的，同样显然在天主身上也是必要的，因为人从祂那里得到榜样。因为如同我们应当约束那些在我们权下的人，同样天主也应当约束所有人的过犯。为了做到这一点，祂必须愤怒；因为一个好人在看到别人的过错时自然地会感动和激动。因此他们本应给出这样的定义：愤怒是心灵为约束过错而激发的情感。因为\Person{西塞罗}给出的定义，即\Quote{愤怒是报复的欲望}，与前述定义并无太大差别。但是那种我们可以称之为狂怒或暴怒的愤怒，即使在人类中也不应存在，因为它是完全邪恶的；但那种与纠正恶习有关的愤怒，既不应从人身上去除，也不能从天主身上去除，因为它既有助于人类的事务，也是必要的。\par

% structure: p0079 source=h2 target=section reason=relative-heading-rank; base=h2

\section{第十八章——论惩罚过犯不能没有愤怒}

他们说，既然过犯可以在没有愤怒这种情感的情况下被纠正，还需要愤怒做什么呢？但是没有人能平静地看到任何人犯过。这或许对于主持法律的人来说是可能的，因为事情不是在他眼前发生的，而是从别处作为可疑之事呈报给他。也没有任何邪恶如此明显，以至于没有辩护的余地；因此法官可能不会对有可能证明无罪的人动怒；当被察觉的罪行曝光后，他不再运用自己的意见，而是运用法律的意见。可以承认，他做他所做的事而没有愤怒；因为他有可遵循的。而我们毫无疑问，当家里有人犯过时，无论我们看到或察觉到，都必须感到愤慨；因为看到罪本身就是不合宜的。因为那完全不为所动的人，要么是赞同过犯——这更可耻、更不义，要么是避免斥责的麻烦——而平静的精神和安宁的心灵除非被愤怒激发，否则会鄙视并拒绝这种麻烦。但是当一个人被感动，却因不合时宜的宽仁而过于频繁或总是给予赦免时，他显然既毁灭了那些他正在助长其胆大妄为以犯更大罪行的人的生命，也为自己带来了永久的烦恼之源。因此，在罪过的情况下抑制愤怒是有过错的。\par

塔兰托的阿基塔斯受到赞扬，当他发现自己的庄园一切都被毁坏时，斥责管家的过错，说：\Quote{你这可怜虫，若我不是在愤怒中，我早将你打死了。}他们认为这是克制的非凡典范；但受到权威的影响，他们没有看到他的言行是何等愚蠢。因为（如\Person{柏拉图}所说）没有智者会因已有过犯而惩罚，而是为防止过犯发生，因此显而易见这位智者树立了何等恶劣的榜样。因为如果奴隶们察觉到他们的主人在不愤怒时使用暴力，而在愤怒时却克制暴力，那么他们显然不会犯小过，免得挨打；却会犯最大的过，以激起那乖戾之人的愤怒，从而逍遥法外。但是，如果他是在盛怒之下，我应称赞他给了愤怒空间，使内心的激动通过时间的间隔而平静下来，他的惩罚可能被限制在适度的范围内。因此，由于愤怒的强烈程度，不应立即施加惩罚，而应延迟，以免对罪犯造成过度的痛苦，或在惩罚者心中激起狂怒的爆发。但是，一个人因小过而受罚，却因大过而免受惩罚，这哪里公平智慧呢？但是，如果他学习了事物的本性和原因，他绝不会宣称如此不合适的克制，使一个邪恶的奴隶因主人生气而欢喜。因为正如天主赋予人的身体许多不同的感官，为生活的使用所必需，同样祂也赋予灵魂各种情感，使生活的进程得以调节；并且正如祂赋予欲望以生育后代，同样祂赋予愤怒以约束过犯。\par

但他们既然不知善恶之事的结局，就为败坏与享乐而运用感官欲望，同样也为了施加伤害而运用愤怒与激情，对他们视为仇恨的对象发怒。因此，他们甚至对没有冒犯的人、对与自己同等的人、甚至对高于自己的人发怒。于是他们每日冲入恶行，悲剧常常由此而生。因此，如果阿尔基塔斯\Person{阿尔基塔斯}在被任何一位公民或同等之辈伤害而发怒时，能克制自己，并以忍耐缓和怒火的猛烈，那他是值得称赞的。这种自制是光荣的，它阻止了任何迫在眉睫的大恶；但不制止奴隶和孩子的过错却是一种过失；因为他们若不受惩罚而逃脱，便会走向更大的恶。在这种情况下，愤怒不应被抑制；即使它处于不活跃状态，也必须被激发。但我们关于人所说的，也关于天主而言，祂造人相似于自己。我暂且不提天主的形象，因为斯多葛派说天主没有形状，如果我们想要驳斥他们，就会产生另一个重大主题。我只就灵魂而言。如果属于天主的特性是反思、智慧、理解、预见、卓越，而所有动物中只有人拥有这些特性，那么他便是按天主的肖像受造的；但正因如此，他走向恶习，因为他混杂了来自尘土的脆弱，除非他被同一的天主以正义的诫命所浸润，否则他无法保持从天主所领受之物的纯洁与无玷。\par

% structure: p0083 source=h2 target=section reason=relative-heading-rank; base=h2

\section{第十九章——论灵魂与身体，以及天主眷顾。}

但既然人由两部分组成，即灵魂与身体，如我们所说，德行包含在其中一方，恶习包含在另一方，它们彼此对立。因为灵魂的善性——在于抑制欲望——与身体相悖；身体的善性——在于各种享乐——与灵魂为敌。但若灵魂的德行抵挡了欲望并抑制了它们，人便真正相似于天主。由此可见，人的灵魂，既然能拥有神圣的德行，便不是可朽的。但这里有区别：既然德行伴随着苦涩，而享乐的吸引力甘甜，多数人便被征服而倾向于甘美；然而，那些将自己交付给身体和尘世事物的人被压向地面，不能获得天主恩赐的眷顾，因为他们已用恶习的污秽污染了自己。但那些跟随天主、服从天主、轻视身体欲望、以德行胜于享乐、保持纯洁与正义的人，天主视他们为相似于自己。\par

因此，既然祂订立了至圣的法律，并愿所有人纯洁而乐善好施，祂看见祂的法律被藐视、德行被拒绝、享乐成为追求目标时，难道可能不愤怒吗？但如果祂是世界的治理者，正如祂应当是的那样，祂定然不会藐视在整个世界中最重要的事物。如果祂有预见，如天主所应有的那样，显然祂关怀人类的利益，为使我们的人生更丰富、更美好、更安全。如果祂是众人的父和天主，祂无疑因人的德行而喜悦，因人的恶习而被激怒。因此祂爱义人，恨恶人。有人说：不需要恨；因为祂已经一劳永逸地为善人设定了赏报，为恶人设定了惩罚。但若有人生活正义而纯洁，却既不敬拜天主，也不在意祂——如同阿里斯提得斯\Person{阿里斯提得斯}、提蒙\Person{提蒙}及其他哲学家——他是否能因服从了天主的法律却藐视了天主本人而免于惩罚？因此，确有某事可令天主对那如同信赖自己的正直而反抗祂的人发怒。如果祂能因这人的骄傲而向他发怒，那么对于那连同立法者一起藐视法律的罪人，岂不更当如此？法官不能赦免过犯，因为他受他人意志的支配。但天主能赦免，因为祂是自己法律的仲裁者和审判者；当祂订立法律时，祂并未剥夺自己的一切权能，而是拥有赦免的自由。\par

% structure: p0086 source=h2 target=section reason=relative-heading-rank; base=h2

\section{第二十章——论过犯与天主的仁慈。}

如果祂能赦免，那么祂也能发怒。那么，有人会说，为什么常常发生罪人得福、虔诚者受苦的事呢？因为逃亡者和被剥夺继承权的人生活无度，而在父亲或主人管教下的人则生活更严格、更节俭。因为德行通过患难得到证明和坚定，恶习则通过享乐。然而，犯罪者不应指望永不受罚，因为没有持久的幸福。\par

\Quote{但事实上，人总应期待最后的日子，无人应在死亡和最后的葬礼之前被称为有福。}\par

正如那位不失优雅的诗人所说。结局证明幸福，无人能逃脱天主的审判，无论是生前还是死后。因为祂有权能将活人从高处推下，也能以永罚惩罚死人。不，他说：如果天主发怒，祂应该立即施行报复，并按各人的罪过惩罚每人。但（有人回应）如果祂这样做，就没有人能存活。因为没有人不在任何方面犯罪，而且有许多事激发人去犯罪——年龄、放纵、匮乏、机会、报酬。我们所穿戴的肉体之脆弱如此易于犯罪，若非天主宽容此必然性，或许存活之人将少之又少。为此，祂极其忍耐，克制自己的愤怒。因为祂内有圆满的德行，必然祂的忍耐也是圆满的，忍耐本身也是一种德行。有多少人从罪人后来成为义人；从害人者成为好人；从恶人成为节制的人！有多少早年卑劣、为众人评判所定罪的人，后来竟变得值得称赞！但显然，若每次过犯都紧随惩罚，这事就不可能发生。\par

公开的法律定那些明显有罪之人的罪；但有大量的人其过犯是隐藏的，有大量的人以恳求或贿赂遏止控告者，有大量的人凭恩宠或势力逃避司法。但若天主的审判应定罪所有逃脱人罚的人，那么地上的人将所剩无几，甚至无人。简言之，单单那个毁灭人类种族的理由也可能是正义的：人藐视永生的天主，却将神圣的尊荣归于地上易朽的偶像，仿佛它们来自天上，崇拜人手所造的物品。尽管天主，他们的创造者，使他们具有高昂的面容和正直的身形，并将他们提升到对天主的默观和对天主的认识，他们却像牲畜一样，宁愿弯腰向地。因为那人的位置是低下的、弯曲的、向下倾斜的，他转离对天堂和天主、他圣父的注视，却崇拜地上的事物——那些他本应践踏的事物，即由泥土所造所塑之物。因此，在如此巨大的不敬和如此巨大的人类的罪恶中，天主的忍耐达到了这个目的：人谴责自己过去生活的错误，改正自己。总之，有许多人公义而善良；这些人已放弃对地上事物的崇拜，承认唯一真天主的威严。但尽管天主的忍耐非常巨大且极为有益，然而祂虽迟延，终究惩罚有罪之人，当祂看见他们不可救药时，不允许他们再继续下去。\par

% structure: p0091 source=h2 target=section reason=relative-heading-rank; base=h2

\section{第二十一章——论天主的愤怒与人的愤怒}

还剩下一个问题，也是最后一个问题。因为有人或许会说，天主远非发怒，祂甚至在祂的诫命中禁止人发怒。我可以说，人的愤怒应当被克制，因为他常常不义地发怒；而且他立即有情绪，因为他只是暂时的。因此，唯恐那些卑微者、中等地位的人和伟大的君主在愤怒中所做的事发生，他的怒气应当被调节和压制，免得他失去理智，犯下某种不可赎的罪行。但天主发怒不是短暂的，因为祂是永恒的，具有完美的德行，并且祂从不无端发怒。但无论如何，事情并非如此；因为如果祂完全禁止愤怒，那么祂自己就在某种程度上是祂自己作品的批判者，因为祂从起初就将愤怒安置在人的肝脏中，因为人们相信这种情绪的起因包含在胆汁的湿气中。因此，祂并未完全禁止愤怒，因为这种情感是必然赋予的，但祂禁止我们坚持愤怒。因为凡人的愤怒应当是必死的；因为如果它持久，敌意就加剧至永久的毁灭。再者，当祂命令我们发怒，却不要犯罪时，显然祂并未将愤怒连根拔起，而是约束它，使我们在每次纠正中都能保持节制和正义。因此，那命令我们发怒的一位，显然祂自己也是发怒的；那命令我们迅速息怒的一位，显然祂自己是容易息怒的：因为祂所命令的这些事，对于社会的利益是正义和有益的。\par

但是因为我曾说，天主的愤怒并非像人的愤怒那样短暂——人瞬间激动，并因自身的软弱难以轻易自制——我们应当明白，因为天主是永恒的，祂的愤怒也永远存留；但另一方面，因为祂赋有最高卓越性，祂控制自己的愤怒，不被愤怒支配，而是按照祂的意志调节愤怒。显然这与刚才所说并无矛盾。因为如果祂的愤怒完全是不朽的，那么在过犯之后就不会有和解或仁慈的余地，尽管祂亲自命令人在日落前和解。但天主的愤怒永远存留于那些一直犯罪的人。因此，平息天主的愤怒不是靠乳香或牺牲，不是靠昂贵的供品——这些东西都是可朽坏的——而是靠道德的改正：那停止犯罪的人使天主的愤怒成为必死的。为此，祂不立即惩罚每一个有罪的人，使人有机会恢复理智并改正自己。\par

% structure: p0094 source=h2 target=section reason=relative-heading-rank; base=h2

\section{第二十二章——论诸罪以及关于它们的西比尔诗句的吟诵}

我最亲爱的\Person{多纳图斯}，这是我对天主义怒所要说的，好使你知道如何驳斥那些主张天主无情的人。现在只剩下效法西塞罗的做法，以结语收尾。正如他在\WorkTitle{图斯库兰论辩集}中谈论死亡主题时所做的那样，我们在这部著作中也应当提出那可信的神圣见证，以驳斥那些人的主张——他们相信天主没有怒气，因而摧毁了一切宗教；我们已经表明，没有宗教，我们便或与野兽同样野蛮，或与牲畜同样愚昧。因为只有在宗教中——即在对至高天主的认识中——才有智慧。所有先知，充满圣神，所说的无非是天主对义人的恩宠和对不虔敬者的义怒。他们的见证对我们而言固然足够，但那些以头发和服装炫耀智慧的人并不相信，因此有必要以理性和论据来驳斥他们。因为他们行事如此颠倒，竟让人的事物赋予神性事物权威，而神性事物本应赋予人的事物权威。但我们如今且放下这些，免得对他们毫无效果，并且话题无限拉长。因此，让我们寻找那些他们能够相信，或至少不能反对的见证吧。\par

数量众多且具分量的作者都曾提及\OriginalTerm{西比尔}{Sibyls}；希腊人中，有\Person{阿里斯通}（\Place{希俄斯}人）和\Person{阿波罗多罗斯}（\Place{埃律特赖}人）；我们自己的作者中，有\Person{瓦罗}和\Person{费内斯特拉}。他们都记述说，\Place{埃律特赖}的西比尔较其他西比尔更为卓越和高贵。事实上，\Person{阿波罗多罗斯}以其为本国同胞而自豪。\Person{费内斯特拉}也记载说，元老院曾派遣使者前往\Place{埃律特赖}，以便将此西比尔的作品运至\Place{罗马}，并委托执政官\Person{库里奥}与\Person{屋大维}负责将其安放在当时由\Person{昆图斯·卡图卢斯}负责重建的\Place{卡皮托利丘}上。在她的著作中，发现了关于至高天主与世界创造者的如下诗句：——\par

\Quote{那住在天上、不朽且永恒的创造者，向善人施予更丰厚的赏报，却向恶人与不义者激起忿怒与怒气。}\par

又在另一处，当她列举天主特别为之动怒的行为时，她引入了这些话：——\par

\Quote{避免不合法的服务，要事奉永生天主。戒绝奸淫与不洁，养育纯洁的子孙后代；不可杀人；因为那不朽者必向每一个犯罪的人发怒。}\par

因此，祂向罪人发怒。\par

% structure: p0101 source=h2 target=section reason=relative-heading-rank; base=h2

\section{第二十三章——论天主义怒与罪的刑罚，并引述西比尔们的相关诗篇；此外，还有责备与劝勉。}

但由于博学之士记述曾有众多西比尔，仅凭一位的见证，恐不足以证实我们所欲确立的真理。事实上，库迈的西比尔的作品中写有关于罗马人命运的卷册，但被秘密保存；而其他所有西比尔的著作，大部分并不禁止公开传阅。在这些作品中，另一位西比尔因人类的不虔敬而宣告天主义怒攻击万邦，如此起头：——\par

\Quote{既然伟大的义怒将临于悖逆的世界，我向末代宣告天主的诫命，从这城到那城向众人预言。}\par

另一位西比尔也说，洪水是由天主对上古不义之人的义怒引起的，为要灭绝人类的邪恶：——\par

\Quote{自从那天上的天主对城邑和众人发怒，洪水暴发，海洋覆盖大地。}\par

同样地，她预言将来要发生一场大火，用以再次消灭人类的不虔敬：——\par

\Quote{到某一时刻，天主不再平息祂的怒气，反而使之加剧，毁灭人类种族，用火将其全部焚毁。}\par

关于这一点，\Person{奥维德}提到\OriginalTerm{朱庇特}{Jupiter}时说：——\par

\Quote{他也记得，命定之时将至，那时海洋、大地与天宫将被火吞噬，那精巧构造的世界框架将陷于危险。}\par

这必将在至高者所受的尊崇与敬拜在人间消失之时发生。然而，同一位西比尔证实，祂因行为的悔改与自我改善而息怒，又补充了这些话：——\par

\Quote{但是，你们这些有死之人，快发怜悯，回心转意吧，不要使伟大的天主动尽各样的怒气。}\par

稍后又言：——\par

\Quote{祂不会毁灭，反而要再次抑制祂的怒气，只要你们都在心中践行可贵的虔诚。}\par

另一位西比尔宣称，应当爱那为天上与地上万有之父的，免得祂的义怒兴起，使人类遭毁灭：——\par

\Quote{以免那不朽的天主发怒，毁灭整个人类种族，连同他们的生命与无耻的族类，我们理当爱那智慧、永生的天主圣父。}\par

由此可见，那些哲学家的论调是枉然的，他们幻想天主没有怒气，并在祂的其他赞美中加上那最无用的，减损了祂对人事最有益的属性——即威严本身赖以存在的属性。因为人间的王国与政府，若非以恐惧守护，便会倾覆。夺去国王的怒气，他不仅不再被服从，甚至会被从高位上推翻。不但如此，若从任何卑微的人身上夺去这种情感，谁不会抢夺他？谁不会嘲笑他？谁不会以伤害待他？这样，他将既无衣着，也无居所，也无食物，因为别人会夺走他的一切；我们更不能设想，天上的政府之威严能在没有怒气与恐惧的情况下存在。当\Place{米利都}的\Person{阿波罗}被问及犹太人的宗教时，他在答复中加入了这些话：——\par

\Quote{天主，普世的君王与圣父，大地、天空与海洋在祂面前战栗，塔塔罗斯的深渊与恶魔都畏惧祂。}\par

如果祂是如哲学家所主张的那样温和，那么不仅魔鬼和如此大能的仆役，就连天地和整个宇宙体系，又怎会在祂面前战栗呢？因为如果没有人出于强迫而服从另一个人的服务，那么一切统治都存在于恐惧之中，而恐惧存在于愤怒之中。因为如果有人不向不愿服从的人生怒，就不可能强迫其服从。让任何人审视自己的情感；他立刻就会明白，没有愤怒和惩罚，没有人能被征服而服从另一个人的命令。因此，哪里没有愤怒，哪里就没有权威。但天主有权威；因此祂也必有愤怒，权威就在于愤怒。因此，不要让任何人因哲学家空洞的喋喋不休而被诱导，训练自己轻视天主，这是最大的不敬。我们都必须爱祂，因为祂是我们的圣父；敬畏祂，因为祂是我们的主：既因祂慷慨而尊崇祂，又因祂严厉而畏惧祂：祂的每一个品性都值得敬畏。谁能保持他的虔敬，却不爱生命的父？或者谁能不受惩罚地轻视祂，祂作为万有的统治者，拥有真实而永恒的大能？如果你视祂为圣父，祂供给我们进入所享受的光明：藉着祂我们生活，藉着祂我们进入这个世界。如果你默想祂为天主，是祂以无数的资源养育我们：是祂维持我们，我们住在祂的家中，我们是祂的家人；如果我们没有如所当行的那样顺从，没有如我们的主人和圣父无尽的功德所要求的那样勤于职守：然而，我们若保持对祂的敬拜和认识，若放下卑下和属世的事务与财物，默想天上永恒的神圣事物，这对我们获得宽恕大有裨益。为使我们能够这样做，我们必须跟随天主，必须敬拜和爱慕天主；因为在祂内有万物的实体，德行的原理，以及一切美善的根源。\par

有什么在能力上比天主更大，在理性上更完美，在明净上更光明？既然祂生了我们归于智慧，产生了我们归于正义，人就不许可离弃赐予智慧和生命的天主，而去服侍属地的、易朽坏的事物，或者专注于寻求暂时之物，偏离纯真与虔敬。邪恶致命的享乐不能使人幸福；富裕，欲望的煽动者，也不能；虚空的野心也不能；易逝的荣誉也不能，人的灵魂被诱惑、被肉身奴役，被判受永死：唯有纯真与正义，其合法且应得的赏报是不朽，这是天主从起初为圣洁无玷的心灵所指定的，它们使自身纯洁，不受恶习和一切属地的污秽玷染。不能分享这属天永恒赏报的，是那些以暴力、欺诈、劫掠和欺骗的行为玷污了良心的人；以及那些通过对人施加伤害、不敬的行为，在自己身上烙下不可磨灭的污点的人。因此，所有希望配得称为智者、称为人的人，都应当轻视易朽之物，践踏属地的物，鄙视卑贱的事，以便能够与天主结合，享有最幸福的关系。\par

愿不敬与不和被除去；愿动荡致命的纷争被平息，人类社会和公共联盟的神圣结合因这些纷争而被侵入、分裂、分散；尽我们所能，让我们力求成为良善和慷慨的人：如果我们有财富和资源的供应，不要将之用于一个人的享乐，而要施予许多人的福祉。因为享乐如同它所服侍的身体一样短暂。但正义与仁慈如同心灵和灵魂一样不朽，它们藉着善行达到与天主的相似。愿我们奉献天主，不在圣殿，而在我们的心中。凡人手所做的一切都是可朽坏的。让我们洁净这殿，它不是被烟尘败坏，而是被邪念败坏；它不是被燃烧的蜡烛照亮，而是被智慧的光明照亮。并且如果我们相信天主常在这殿里，祂的神性洞悉心中的秘密，我们就要这样生活，以便常使祂仁慈，永不畏惧祂的愤怒。\par

\SourceRecord{Translated by William Fletcher.；Ante-Nicene Fathers；Vol. 7.；Edited by Alexander Roberts, James Donaldson, and A. Cleveland Coxe.；Buffalo, NY: Christian Literature Publishing Co.,；1886.；<http://www.newadvent.org/fathers/0703.htm>.}
